% Shared definitions for study-guide.tex and exercises.tex
% (Week 10).
% Same colors and box look as script.tex; no lecturer cues.

\definecolor{navy}{RGB}{24,48,75}
\definecolor{teal}{RGB}{0,120,125}
\definecolor{gold}{RGB}{190,130,20}

\addtokomafont{disposition}{\color{navy}}
\setkomafont{section}{\Large\bfseries\color{navy}}

\newcommand{\Om}{\Omega}
% Expectation, variance and standard deviation, as in Ross: E[X], Var(X), SD(X)
\newcommand{\E}{\operatorname{E}}
\newcommand{\Var}{\operatorname{Var}}
\newcommand{\SD}{\operatorname{SD}}
\newcommand{\myemail}{tiborsandor2001@gmail.com}

% Exercise statement (identical in both versions)
\newtcolorbox{task}[1]{enhanced,breakable,colback=navy!4,colframe=navy,
	fonttitle=\bfseries,title={#1}}
% Finance example
\newtcolorbox{fin}[1]{enhanced,breakable,colback=teal!5,colframe=teal,
	fonttitle=\bfseries,title={#1}}
% Key formulas at the start of each part
\newtcolorbox{key}[1]{enhanced,breakable,colback=white,colframe=navy!70,
	fonttitle=\bfseries,title={#1}}
% Final answer (solutions version only)
\newtcolorbox{answer}{enhanced,colback=teal!6,colframe=teal,boxrule=.5pt,
	left=4pt,right=4pt,top=2pt,bottom=2pt}
% Remarks: alternative methods, interpretation
\newtcolorbox{remark}{enhanced,breakable,colback=gold!8,colframe=gold,
	boxrule=.4pt,left=4pt,right=4pt,top=2pt,bottom=2pt,fontupper=\small}

% Labels inside solutions
\newcommand{\soltitle}{\par\textbf{\color{teal}Solution.}\par\nobreak}
\newcommand{\idea}{\textbf{Idea.}\ }
\newcommand{\step}[1]{\par\textbf{#1}\ }

% Small Venn diagram of two events E, F in Omega; #1 = shading commands
\newcommand{\venn}[1]{%
	\begin{tikzpicture}[scale=.5,baseline=(O.base)]
		\node (O) at (.5,0) {\phantom{x}};
		#1
		\draw[thick] (-1.5,-1.15) rectangle (2.5,1.15);
		\draw[thick,navy] (0,0) circle (.85);
		\draw[thick,navy] (1,0) circle (.85);
		\node at (-.35,0) {\scriptsize $E$};
		\node at (1.35,0) {\scriptsize $F$};
		\node at (2.15,.8) {\scriptsize $\Omega$};
	\end{tikzpicture}}
\newcommand{\shadeall}[1]{\fill[#1] (-1.5,-1.15) rectangle (2.5,1.15);}

% Pictures of continuous distributions (teal areas, navy curves, gold marks).
% The same definitions are in beamer.tex and preamble.tex.
% A density f(x) = x e^{-x} (x >= 0) and the area P(a <= X <= b) under it,
% a = 1.5, b = 3.2 (area 0.387).
\newcommand{\figdensity}{%
  \begin{tikzpicture}[x=.85cm,y=5cm]
    \fill[teal!25] (1.5,0) -- plot[domain=1.5:3.2,samples=40]
      (\x,{\x*exp(-\x)}) -- (3.2,0) -- cycle;
    \draw[->] (-.2,0)--(6.4,0) node[right,font=\scriptsize] {$x$};
    \draw[->] (0,-.02)--(0,.47);
    \draw[very thick,navy] (-.2,0)--(0,0) plot[domain=0:6.1,samples=80]
      (\x,{\x*exp(-\x)});
    \node[navy,font=\scriptsize] at (1.7,.42) {$f(x)$};
    \node[below,font=\scriptsize] at (1.5,0) {$a$};
    \node[below,font=\scriptsize] at (3.2,0) {$b$};
    \draw[teal,->] (4.1,.24)--(2.5,.13);
    \node[right,font=\scriptsize,align=left] at (3.9,.29)
      {area $=$\\$P(a\le X\le b)$};
  \end{tikzpicture}}
% The distribution function F(x) = 1 - (1+x) e^{-x} of the same density;
% the arrow marks F(b) - F(a), the same area.
\newcommand{\figcdf}{%
  \begin{tikzpicture}[x=.8cm,y=1.9cm]
    \draw[->] (-.2,0)--(6.4,0) node[right,font=\scriptsize] {$x$};
    \draw[->] (0,-.05)--(0,1.2);
    \draw[gray!50,dashed] (0,1)--(6.1,1);
    \node[left,font=\scriptsize] at (0,1) {$1$};
    \draw[very thick,navy] (-.2,0)--(0,0) plot[domain=0:6.1,samples=80]
      (\x,{1-(1+\x)*exp(-\x)});
    \node[navy,font=\scriptsize] at (1.0,.62) {$F(x)$};
    \draw[gold,densely dotted] (1.5,0)--(1.5,0.442)--(6.1,0.442)
      (3.2,0)--(3.2,0.829)--(6.1,0.829);
    \node[below,font=\scriptsize] at (1.5,0) {$a$};
    \node[below,font=\scriptsize] at (3.2,0) {$b$};
    \draw[gold,very thick,<->] (5.6,0.442)--(5.6,0.829);
    \node[right,font=\scriptsize] at (5.6,.635) {$F(b)-F(a)$};
  \end{tikzpicture}}
% Uniform density on [alpha, beta]; the shaded interval [c, d] has
% probability (d - c)/(beta - alpha); the gold mark is E[X].
\newcommand{\figuniform}{%
  \begin{tikzpicture}[x=.9cm,y=1.6cm]
    \fill[teal!25] (2.1,0) rectangle (3.1,.8);
    \draw[->] (-.2,0)--(5.4,0) node[right,font=\scriptsize] {$x$};
    \draw[->] (0,-.05)--(0,1.25);
    \draw[very thick,navy] (-.2,0)--(1,0) (1,.8)--(4.6,.8) (4.6,0)--(5.2,0);
    \draw[navy,densely dotted] (1,0)--(1,.8) (4.6,0)--(4.6,.8);
    \draw[gray!50,dashed] (0,.8)--(1,.8);
    \node[left,font=\scriptsize] at (0,.8) {$\frac{1}{\beta-\alpha}$};
    \node[below,font=\scriptsize] at (1,0) {$\alpha$};
    \node[below,font=\scriptsize] at (4.6,0) {$\beta$};
    \node[below,font=\scriptsize] at (2.1,0) {$c$};
    \node[below,font=\scriptsize] at (3.1,0) {$d$};
    \node[font=\tiny] at (2.6,.4) {$\frac{d-c}{\beta-\alpha}$};
    \draw[gold,very thick,dashed] (2.8,.8)--(2.8,1.12);
    \node[above,font=\tiny] at (2.8,1.1) {$\E[X]=\frac{\alpha+\beta}{2}$};
  \end{tikzpicture}}
% Exponential density with rate lambda (drawn with lambda = 1): the shaded
% tail beyond x has area e^{-lambda x}; the gold mark is E[X] = 1/lambda.
\newcommand{\figexponential}{%
  \begin{tikzpicture}[x=.9cm,y=1.5cm]
    \fill[teal!25] (1.6,0) -- plot[domain=1.6:5,samples=40]
      (\x,{exp(-\x)}) -- (5,0) -- cycle;
    \draw[->] (-.2,0)--(5.4,0) node[right,font=\scriptsize] {$t$};
    \draw[->] (0,-.05)--(0,1.25);
    \draw[very thick,navy] (-.2,0)--(0,0);
    \draw[very thick,navy] plot[domain=0:5,samples=60] (\x,{exp(-\x)});
    \node[left,font=\scriptsize] at (0,1) {$\lambda$};
    \node[navy,font=\scriptsize] at (1.05,.75) {$\lambda e^{-\lambda t}$};
    \draw[gold,very thick,dashed] (1,0)--(1,.37);
    \node[below,font=\scriptsize] at (1,0) {$\frac1\lambda$};
    \node[below,font=\scriptsize] at (1.6,0) {$x$};
    \draw[teal,->] (3.3,.55)--(2.3,.08);
    \node[above,font=\scriptsize] at (3.6,.52) {$P(X>x)=e^{-\lambda x}$};
  \end{tikzpicture}}
% Normal density with mean mu and standard deviation sigma: the bands
% mu +- sigma, mu +- 2 sigma, mu +- 3 sigma and their probabilities.
\newcommand{\fignormal}{%
  \begin{tikzpicture}[x=.8cm,y=5.2cm]
    \fill[teal!10] (-3,0) -- plot[domain=-3:3,samples=80]
      (\x,{exp(-\x*\x/2)/sqrt(2*pi)}) -- (3,0) -- cycle;
    \fill[teal!22] (-2,0) -- plot[domain=-2:2,samples=60]
      (\x,{exp(-\x*\x/2)/sqrt(2*pi)}) -- (2,0) -- cycle;
    \fill[teal!40] (-1,0) -- plot[domain=-1:1,samples=40]
      (\x,{exp(-\x*\x/2)/sqrt(2*pi)}) -- (1,0) -- cycle;
    \draw[->] (-3.7,0)--(3.9,0) node[right,font=\scriptsize] {$x$};
    \draw[very thick,navy] plot[domain=-3.6:3.6,samples=100]
      (\x,{exp(-\x*\x/2)/sqrt(2*pi)});
    \draw[gold,very thick,dashed] (0,0)--(0,.44);
    \node[above,font=\scriptsize] at (0,.44) {$\mu$};
    \foreach \k/\l in {-3/\mu{-}3\sigma,-2/\mu{-}2\sigma,-1/\mu{-}\sigma,
        1/\mu{+}\sigma,2/\mu{+}2\sigma,3/\mu{+}3\sigma} {
        \draw (\k,0)--(\k,-.015);
        \node[below,font=\tiny,inner sep=1pt] at (\k,-.01) {$\l$};
      }
    \draw[navy,<->] (-1,-.1)--(1,-.1);
    \node[fill=white,inner sep=1pt,font=\tiny] at (0,-.1) {$68\%$};
    \draw[navy,<->] (-2,-.145)--(2,-.145);
    \node[fill=white,inner sep=1pt,font=\tiny] at (0,-.145) {$95\%$};
    \draw[navy,<->] (-3,-.19)--(3,-.19);
    \node[fill=white,inner sep=1pt,font=\tiny] at (0,-.19) {$99.7\%$};
  \end{tikzpicture}}
% The standard normal Phi(z) as the area left of z, and the symmetry
% Phi(-z) = 1 - Phi(z): the two gold tails have the same area.
\newcommand{\figphi}{%
  \begin{tikzpicture}[x=.3cm,y=3.2cm]
    \fill[teal!30] (-3.3,0) -- plot[domain=-3.3:1,samples=60]
      (\x,{exp(-\x*\x/2)/sqrt(2*pi)}) -- (1,0) -- cycle;
    \draw[->] (-3.5,0)--(3.6,0);
    \draw[very thick,navy] plot[domain=-3.3:3.3,samples=80]
      (\x,{exp(-\x*\x/2)/sqrt(2*pi)});
    \node[below,font=\scriptsize] at (1,0) {$z$};
    \node[below,font=\scriptsize] at (0,0) {$0$};
    \node[font=\tiny] at (-.2,.11) {$\Phi(z)$};
    \node[font=\tiny] at (0,.5) {area left of $z$};
    \begin{scope}[xshift=2.45cm]
      \fill[gold!50] (-3.3,0) -- plot[domain=-3.3:-1,samples=40]
        (\x,{exp(-\x*\x/2)/sqrt(2*pi)}) -- (-1,0) -- cycle;
      \fill[gold!50] (1,0) -- plot[domain=1:3.3,samples=40]
        (\x,{exp(-\x*\x/2)/sqrt(2*pi)}) -- (3.3,0) -- cycle;
      \draw[->] (-3.5,0)--(3.6,0);
      \draw[very thick,navy] plot[domain=-3.3:3.3,samples=80]
        (\x,{exp(-\x*\x/2)/sqrt(2*pi)});
      \node[below,font=\scriptsize] at (-1,0) {$-z$};
      \node[below,font=\scriptsize] at (1,0) {$z$};
      \node[font=\tiny] at (0,.5) {$\Phi(-z)=1-\Phi(z)$};
    \end{scope}
  \end{tikzpicture}}
% The unit square: the joint density of two variables on [0,1]^2 and a
% region A (below the diagonal) whose probability is a double integral.
\newcommand{\figjointsquare}{%
  \begin{tikzpicture}[scale=1.5]
    \fill[teal!12] (0,0) rectangle (1.4,1.4);
    \fill[teal!45] (0,0) -- (1.4,0) -- (1.4,1.4) -- cycle;
    \draw[->] (-.1,0) -- (1.7,0) node[right,font=\scriptsize] {$x$};
    \draw[->] (0,-.1) -- (0,1.7) node[above,font=\scriptsize] {$y$};
    \draw[thick,navy] (0,0) rectangle (1.4,1.4);
    \draw[navy,densely dashed] (0,0) -- (1.4,1.4);
    \node[below,font=\scriptsize] at (1.4,0) {$1$};
    \node[left,font=\scriptsize] at (0,1.4) {$1$};
    \node[font=\scriptsize] at (1.0,.35) {$A$};
    \node[font=\tiny,rotate=45] at (.55,.72) {$y=x$};
  \end{tikzpicture}}

\title{Mathematics for Finance}
\author{Tibor SÁNDOR\\[2pt]
	\normalsize Corvinus University of Budapest\\
	\normalsize\href{mailto:\myemail}{\texttt{\myemail}}}

% The exercises (one file each in exercises/) and the commands that place them
% Exercises of Week 10, shared by beamer.tex, the study guide, the script and
% exercises.tex. Every exercise lives in a file of its own,
% exercises/NN-slug.tex, numbered 10.NN (see exercises.lua for the docstring).
% The body of an exercise file:
%
%   \begin{problem} ... \end{problem}      the statement, used everywhere;
%                                          inside it, \begin{context} ...
%                                          \end{context} is background that
%                                          the slides leave out
%   \begin{slideproblem} ... \end{slideproblem}
%                                          (optional) a shorter statement
%                                          for the slides
%   \begin{slides} ... \end{slides}        one frame of the deck: \exstatement
%                                          places the statement; repeat the
%                                          environment for more frames
%   \begin{solution} ... \end{solution}    the detailed solution (study
%                                          guide, exercise PDFs, website)
%
% Commands for the documents (the argument is the slug: 01-license-plates.tex
% has the slug license-plates):
%
%   \exref{slug}                 the number, e.g. 10.1, for "Exercise~\exref{...}"
%   \extitle{slug}               the title from the docstring
%   \exerciseframe[opts]{slug}   beamer: the frames of the exercise
%   \studyexercise{slug}         study guide: heading, statement, solution
%   \scriptexercise{slug}        script: the statement box, as on the slide
%   \allexercises                exercises.tex: every exercise, in order; on
%                                the website statement and solution with
%                                markers where the pipeline splits the HTML
%                                (\collectionexercise), in the PDF plain:
%                                each statement, followed by its solution if
%                                \ifsolutions (exercises-solutions.tex)

\def\exweek{10}

\directlua{exercises = dofile("exercises/exercises.lua")
  exercises.scan("exercises", "\exweek")}

\ExplSyntaxOn

\msg_new:nnn { exercises } { unknown } { Unknown~exercise~`#1'. }

\cs_new_protected:Npn \__ex_missing:n #1
  { \msg_error:nnn { exercises } { unknown } {#1} ?? }
\cs_new:Npn \exref #1
  {
    \cs_if_exist:cTF { exnum@#1 }
      { \use:c { exnum@#1 } } { \__ex_missing:n {#1} }
  }
\cs_new:Npn \extitle #1
  {
    \cs_if_exist:cTF { extitle@#1 }
      { \use:c { extitle@#1 } } { \__ex_missing:n {#1} }
  }

% ---------------------------------------------------------------------------
% Loading an exercise: its parts are stored, nothing is typeset. The slides
% typeset the stored parts; the other documents read the part they need from
% the file again (\__ex_input_part:nn), as stored tables do not survive lwarp.

\tl_new:N \g__ex_current_tl
\tl_new:N \g__ex_problem_tl
\tl_new:N \g__ex_slideproblem_tl
\tl_new:N \g__ex_solution_tl
\seq_new:N \g__ex_slides_seq
\box_new:N \l__ex_scratch_box

\cs_new_protected:Npn \__ex_load:n #1
  {
    \cs_if_exist:cF { exfile@#1 } { \__ex_missing:n {#1} }
    \tl_gset:Nn \g__ex_current_tl {#1}
    \tl_gclear:N \g__ex_problem_tl
    \tl_gclear:N \g__ex_slideproblem_tl
    \tl_gclear:N \g__ex_solution_tl
    \seq_gclear:N \g__ex_slides_seq
    \vbox_set:Nn \l__ex_scratch_box
      {
        \DeclareDocumentEnvironment { problem } { +b }
          { \tl_gset:Nn \g__ex_problem_tl {##1} } { }
        \DeclareDocumentEnvironment { slideproblem } { +b }
          { \tl_gset:Nn \g__ex_slideproblem_tl {##1} } { }
        \DeclareDocumentEnvironment { slides } { +b }
          { \seq_gput_right:Nn \g__ex_slides_seq {##1} } { }
        \DeclareDocumentEnvironment { solution } { +b }
          { \tl_gset:Nn \g__ex_solution_tl {##1} } { }
        \exp_args:Nv \input { exfile@#1 }
      }
  }

% Typesets one part (problem, slideproblem, solution) of an exercise straight
% from its file; the other parts are skipped.
\cs_new_protected:Npn \__ex_input_part:nn #1#2
  {
    \group_begin:
      \clist_map_inline:nn { problem, slideproblem, slides, solution }
        {
          \str_if_eq:nnTF {##1} {#2}
            { \DeclareDocumentEnvironment {##1} { } { } { } }
            { \DeclareDocumentEnvironment {##1} { +b } { } { } }
        }
      \exp_args:Nv \input { exfile@#1 }
    \group_end:
  }

% Exercise 10.1, or Exercise 10.2 · Finance
\cs_new:Npn \__ex_heading:
  {
    Exercise~\exref { \g__ex_current_tl }
    \cs_if_exist:cT { exfinance@ \g__ex_current_tl } { ~·~Finance }
  }
\prg_new_conditional:Npnn \__ex_if_finance: { T, F, TF }
  {
    \cs_if_exist:cTF { exfinance@ \g__ex_current_tl }
      { \prg_return_true: } { \prg_return_false: }
  }

% ---------------------------------------------------------------------------
% Beamer

% The finance context is printed in the study guide, the exercise PDFs and on
% the website; the
% slides and the script (which shows what the slide shows) leave it out.
\bool_new:N \l__ex_context_bool
\bool_set_true:N \l__ex_context_bool
\DeclareDocumentEnvironment { context } { }
  {
    \bool_if:NTF \l__ex_context_bool
      { \par \textit { Finance~context. } ~ }
      { \__ex_gobble_context:w }
  }
  { \bool_if:NT \l__ex_context_bool { \par } }
\cs_new_protected:Npn \__ex_gobble_context:w #1 \end #2
  { \str_if_eq:nnTF {#2} { context } { \end { context } } { \__ex_gobble_context:w } }

\IfClassLoadedTF { beamer }
  {
    \bool_set_false:N \l__ex_context_bool

    \cs_new:Npn \__ex_block: { \__ex_if_finance:TF { exampleblock } { block } }
    \NewDocumentCommand \exstatement { }
      {
        \begin { \__ex_block: } { \__ex_heading: }
          \tl_if_empty:NTF \g__ex_slideproblem_tl
            { \g__ex_problem_tl } { \g__ex_slideproblem_tl }
        \end { \__ex_block: }
      }

    \tl_new:N \l__ex_frames_tl
    \NewDocumentCommand \exerciseframe { O{t} m }
      {
        \__ex_load:n {#2}
        \tl_clear:N \l__ex_frames_tl
        \seq_map_indexed_inline:Nn \g__ex_slides_seq
          {
            \tl_put_right:Nn \l__ex_frames_tl { \begin{frame}[#1] }
            \tl_put_right:Ne \l__ex_frames_tl
              {
                {
                  \exp_not:N \extitle {#2}
                  \int_compare:nNnT {##1} > { 1 } { \c_space_tl (cont.) }
                }
              }
            \tl_put_right:Nn \l__ex_frames_tl { ##2 \end{frame} }
          }
        \tl_use:N \l__ex_frames_tl
      }
  }
  {
    \cs_new:Npn \__ex_box: { \__ex_if_finance:TF { fin } { task } }
    \cs_new_protected:Npn \__ex_statement_box:n #1
      {
        \begin { \__ex_box: } { \__ex_heading: }
          \__ex_input_part:nn { \g__ex_current_tl } {#1}
        \end { \__ex_box: }
      }

    \NewDocumentCommand \studyexercise { m }
      {
        \__ex_load:n {#1}
        \subsection* { Exercise~\exref{#1}~---~\extitle{#1} }
        \addcontentsline { toc } { subsection }
          { Exercise~\exref{#1}~---~\extitle{#1} }
        \__ex_statement_box:n { problem }
        \soltitle \__ex_input_part:nn {#1} { solution }
      }

    \NewDocumentCommand \scriptexercise { m }
      {
        \__ex_load:n {#1}
        \group_begin:
          \bool_set_false:N \l__ex_context_bool
          \tl_if_empty:NTF \g__ex_slideproblem_tl
            { \__ex_statement_box:n { problem } }
            { \__ex_statement_box:n { slideproblem } }
        \group_end:
      }

    % The website (lwarp): statement and solution of each exercise, between
    % markers where the pipeline splits the HTML.
    \NewDocumentCommand \collectionexercise { m }
      {
        \__ex_load:n {#1}
        \par [[exercise:#1:problem]] \par
        \__ex_input_part:nn {#1} { problem }
        \par [[exercise:#1:solution]] \par
        \__ex_input_part:nn {#1} { solution }
        \par [[exercise:#1:end]] \par
      }
    % The PDF: plain, each statement, followed by its solution in the
    % version with solutions (\ifsolutions). The statement is set off like a
    % quote in Markdown, by a bar on the left (quote=<color>, also used by
    % exercises.tex) in the color of its box in the study guide: navy, or teal
    % for a finance exercise.
    \tcbset
      {
        quote/.style =
          {
            enhanced, breakable, frame~hidden, interior~hidden,
            sharp~corners, boxrule = 0pt, borderline~west = {2pt}{0pt}{#1},
            left = 8pt, right = 0pt, top = 1pt, bottom = 1pt,
            before~skip = 6pt, after~skip = 6pt,
            % the document's paragraph spacing (parskip=half)
            before~upper = { \setlength { \parskip } { .5\baselineskip } }
          }
      }
    \NewDocumentCommand \plainexercise { m }
      {
        \__ex_load:n {#1}
        \subsection* { Exercise~\exref{#1}~---~\extitle{#1} }
        \addcontentsline { toc } { subsection }
          { Exercise~\exref{#1}~---~\extitle{#1} }
        \__ex_if_finance:TF
          { \begin{tcolorbox}[quote=teal] }
          { \begin{tcolorbox}[quote=navy] }
          \__ex_input_part:nn {#1} { problem }
        \end{tcolorbox}
        \legacy_if:nT { solutions }
          {
            \soltitle
            \__ex_input_part:nn {#1} { solution }
          }
      }
    \NewDocumentCommand \allexercises { }
      {
        \IfPackageLoadedTF { lwarp }
          { \directlua { exercises.each("collectionexercise") } }
          { \directlua { exercises.each("plainexercise") } }
      }
  }

\ExplSyntaxOff


% PDF properties and the date on the title page (\documentinfo)
% The properties of every PDF of Week 10: author (with email address), title,
% subject, keywords and the date the document last changed, in the PDF
% (hyperref) and on its title page (\date), and the header and footer of the
% printed documents (below). Each document loads this file after hyperref and
% then calls \documentinfo{<what it is>}, e.g. \documentinfo{Slides}.
%
% The date: the website build (scripts/build-content.mjs) puts
% \def\lastmodified{YYYY-MM-DD} in front of the document, the date of the last
% commit that changed its sources. Compiled locally, it is today.

\newcommand{\courseinfo}{Mathematics for Finance}
\newcommand{\weeklabel}{Week 10}
\newcommand{\weekinfo}{\weeklabel: Continuous random variables, joint distributions, independence}
\newcommand{\authorname}{Tibor SÁNDOR}
\providecommand{\myemail}{tiborsandor2001@gmail.com}

\makeatletter
\ifdefined\lastmodified\else
	\edef\lastmodified{\the\year-\two@digits{\the\month}-\two@digits{\the\day}}
\fi
\def\di@split#1-#2-#3\@nil{%
	\def\di@year{#1}\def\di@month{#2}\def\di@day{#3}}
\expandafter\di@split\lastmodified\@nil
% November 11, 2026
\edef\lastmodifiedtext{%
	\ifcase\numexpr\di@month\relax\or January\or February\or March\or April\or
		May\or June\or July\or August\or September\or October\or November\or
		December\fi\space\number\di@day, \di@year}
% D:20261111000000, the date format of PDF metadata
\edef\di@pdfdate{D:\di@year\di@month\di@day000000}

\newcommand{\documentinfo}[1]{%
	\hypersetup{
		pdftitle={\courseinfo, \weekinfo\space -- #1},
		pdfauthor={\authorname\space<\myemail>},
		pdfsubject={\courseinfo, \weekinfo},
		pdfkeywords={probability, random variables, mathematics for finance,
				Corvinus University of Budapest},
		pdfcreationdate={\di@pdfdate},
		pdfmoddate={\di@pdfdate},
		pdflang={en-US},
	}%
	\date{Last updated: \lastmodifiedtext}%
	\gdef\documentkind{#1}%
}

% Header and footer of the printed documents (KOMA-Script, scrlayer-scrpage;
% the slides have beamer's footline, the website neither): the course, the
% week and the document on the left of the header (the week's topic is on the
% title page), the university on the right; in the footer the author and the
% page as <current>/<total>. The title page gets only the footer.
\newcommand{\documentkind}{}
\IfClassLoadedTF{beamer}{}{\IfPackageLoadedTF{lwarp}{}{%
	\usepackage{lastpage}%
	\usepackage[headsepline=1.5pt,footsepline=1.5pt,plainfootsepline]
		{scrlayer-scrpage}%
	\clearpairofpagestyles
	\setkomafont{pageheadfoot}{\small\sffamily\color{navy}}%
	\setkomafont{pagenumber}{}%
	\setkomafont{headsepline}{\color{navy!35}}%
	\setkomafont{footsepline}{\color{navy!35}}%
	\ihead{\courseinfo\ \textperiodcentered\ \weeklabel\ \textperiodcentered\ 
		\documentkind}%
	\ohead{Corvinus University of Budapest}%
	\ifoot*{\authorname\ \textperiodcentered\ \myemail}%
	\ofoot*{\pagemark/\pageref*{LastPage}}%
	\pagestyle{scrheadings}%
	% Extra stretch for breaking lines with long inline formulas.
	\setlength{\emergencystretch}{1.5em}%
}}
\makeatother

